
\chapter{Установка и первый запуск}
\label{chap:getting-started}

\section{Требования}

Для работы с \CIAO{} необходим Python~3.11 или новее.
Графические примеры используют Tkinter, который обычно
входит в стандартную поставку Python.

Для экспорта диаграмм с улучшенной компоновкой
дополнительно потребуется Node.js. Его установка
рассмотрена ниже.

\section{Установка}

Скачайте проект из репозитория CIAO:

\begin{lstlisting}[language={},numbers=none]
git clone https://github.com/AlexSoulJk/SPbSTU-26-27-CIAO-for-DMLabs.git
cd SPbSTU-26-27-CIAO-for-DMLabs
\end{lstlisting}

Откройте PowerShell в папке проекта и создайте
виртуальное окружение:

\begin{lstlisting}[language={},numbers=none]
py -3.11 -m venv .venv
.\.venv\Scripts\Activate.ps1
\end{lstlisting}

Установите фреймворк:

\begin{lstlisting}[language={},numbers=none]
python -m pip install -e .
\end{lstlisting}

Флаг \code{-e} позволяет изменять исходный код проекта
без повторной установки пакета.

Проверить установку можно запуском тестов:

\begin{lstlisting}[language={},numbers=none]
python -m unittest discover -s tests -v
\end{lstlisting}

\begin{resultbox}
Если тесты завершились без ошибок, можно переходить
к запуску примеров.
\end{resultbox}

\section{Запуск примеров}

В репозитории находятся готовые приложения,
демонстрирующие возможности фреймворка.

Запустите пример вычисления суммы массива:

\begin{lstlisting}[language={},numbers=none]
python -m examples.array_sum.interactive_run
\end{lstlisting}

Или калькулятор с пошаговой визуализацией вычислений:

\begin{lstlisting}[language={},numbers=none]
python -m examples.calculator.interactive_run
\end{lstlisting}

Попробуйте управлять приложениями с помощью кнопок
и проследить, как меняется их состояние.
В следующих главах мы разберём, как они устроены.

\section{Экспорт диаграмм в draw.io}

\CIAO{} позволяет экспортировать автоматные программы
в формат \code{.drawio}. Полученные диаграммы можно
открывать и редактировать в \url{https://app.diagrams.net/}.

Для компоновки диаграмм используется алгоритм ELK
(Eclipse Layout Kernel), поэтому перед первым экспортом
необходимо установить Node.js с \url{https://nodejs.org/}.

После установки Node.js выполните в корне проекта:

\begin{lstlisting}[language={},numbers=none]
npm install
\end{lstlisting}

Для экспорта диаграмм готовых примеров запустите:

\begin{lstlisting}[language={},numbers=none]
python examples\export_drawio.py
\end{lstlisting}

Готовые файлы появятся в каталоге
\RepoPath{examples/schemas/generated_drawio/}.
Их можно открыть в draw.io для просмотра или редактирования.
